This section presents the proposed hybrid model: an overview of its three phases,
the architecture and training of each, the mechanism that distinguishes the
fingerprint encoder from conventional fine tuning, and the design of the ablation
used to isolate each component.

\subsection{Framework Overview}\label{sec:framework}

Figure~\ref{fig:methodology} summarises the proposed method in two steps. Step~1
(top) builds the hybrid in three phases: it pretrains on past LFP cells, then, for
each new cell, forms an early-life fingerprint and fine tunes. Step~2 (bottom)
shows the three outputs this supports, one per deployment question (how healthy
the cell is now, which cells will fail, and how long each cell will last), and
pairs each output with the experiment that validates it. Every battery dataset fed
into the model is marked with a battery icon. The core is the three-phase hybrid
that predicts next-cycle SOH, detailed below and evaluated in
Section~\ref{sec:results}.
Phase~1 pretrains the hybrid network, the convolutional fingerprint encoder
together with the LSTM, end-to-end on the chemistry matched MIT dataset to
acquire general LFP degradation knowledge.
Phase~2 applies the pretrained encoder to the first 30 discharge cycles of a
target cell, producing a fingerprint that initialises the LSTM hidden state and
conditions all subsequent predictions.
Phase~3 fine tunes the LSTM on the target cell's training split and performs
MC-Dropout inference to deliver uncertainty-aware predictions whose
prediction intervals are empirically validated rather than merely assumed.

The same early-life representation supports two capabilities developed later in
the paper. An early screen for manufacturing-defective cells reads the discharge
voltage curve rather than the capacity trajectory (Section~\ref{sec:screening}),
and a remaining-useful-life forecaster replaces the one-step head with a direct
multi-horizon decoder (Section~\ref{sec:rul}). The two transfer levers are
distinct: the base hybrid is pretrained on the chemistry-aligned MIT LFP set,
whereas the RUL decoder is pretrained on a 155-cell LFP foundation corpus drawn
from the MIT, HUST, and SNL collections. These two levers are matched to two
different tasks rather than merged into one: one-step SOH is a local prediction
whose accuracy saturates on these smooth large-format trajectories
(Section~\ref{sec:results}), so the chemistry-aligned MIT corpus is sufficient
and additional data brings no measurable gain, whereas remaining-useful-life
forecasting is a long-horizon, data-limited task whose error falls only when the
corpus is enlarged to the full 155-cell collection
(Section~\ref{sec:rul_foundation}). Because that foundation set contains the MIT
cells used for the base hybrid, the two corpora are nested rather than
independent, so the RUL stage extends the SOH stage rather than competing with
it; we keep the two pretraining stages separate only so that each transfer lever
can be isolated in the ablation.

\begin{figure*}[htbp]
\centering
\includegraphics[width=\textwidth]{methodology.png}
\caption{The proposed method: the three-phase hybrid and its three validated outputs.}
\label{fig:methodology}
\end{figure*}

\subsection{Hybrid Model Architecture}\label{sec:hybrid_arch}

The hybrid model operationalises the two capabilities identified in
Section~\ref{sec:motivation} through three phases, drawn as the Step~1 sequence at
the top of Figure~\ref{fig:methodology}:

\subsubsection{Phase 1: MIT Pretraining}
The complete hybrid architecture (the convolutional fingerprint encoder of
Phase~2 together with the LSTM) is pretrained \textit{end-to-end} on the
40-cell MIT LFP corpus (Section~\ref{sec:datasets}).
For each cell, the first 30 cycles form the fingerprint input and the network
is trained to predict every \textit{subsequent} cycle's SOH from a recent
look-back window, minimising next-cycle mean squared error
(LSTM: $H=64$, $W=20$; Adam, learning rate $10^{-3}$; early stopping on a 10\%
validation split).
This single objective jointly learns two things: general LFP degradation
dynamics in the LSTM, and the mapping from early-life trajectory shape to the
LSTM initial state in the fingerprint encoder.
The encoder weights obtained here are reused, and frozen, in Phase~3.

\subsubsection{Phase 2: Convolutional Fingerprint Encoder}
A 1-D convolutional network (two Conv1D layers with 32 and 16 filters
respectively, kernel size 3, followed by global average pooling and a dense
projection) encodes the first 30 SOH cycles of the
target cell into a 16-dimensional embedding vector.
This encoder is the same one trained jointly in Phase~1; its parameters are
not re-learned per target cell.
Two dense layers project this embedding to the LSTM initial hidden state
$h_0$ and cell state $c_0$, conditioning predictions on the individual
cell's observed early behaviour.

Formally, let $\mathbf{s} = [s_1, s_2, \ldots, s_{30}]^{\top} \in \mathbb{R}^{30}$
denote the first 30 normalised discharge capacities of the target cell.
The encoder computes:
\begin{equation}
  \mathbf{z} = \mathrm{GlobalAvgPool}\!\left(
    \mathrm{Conv}_2\!\left(\mathrm{Conv}_1(\mathbf{s})\right)
  \right) \in \mathbb{R}^{16}
  \label{eq:fingerprint}
\end{equation}
\begin{equation}
  h_0 = \tanh(\mathbf{W}_h\,\mathbf{z} + \mathbf{b}_h), \qquad
  c_0 = \tanh(\mathbf{W}_c\,\mathbf{z} + \mathbf{b}_c)
  \label{eq:h0c0}
\end{equation}
where $\mathbf{z}$ is the fingerprint embedding from Equation~\eqref{eq:fingerprint},
$\mathbf{W}_h$ and $\mathbf{W}_c$ are the weight matrices of the two dense
projection layers, $\mathbf{b}_h$ and $\mathbf{b}_c$ are their bias vectors, and
$h_0$ and $c_0$ are the resulting LSTM initial hidden and cell states.
The critical point is that $\mathbf{s}$ is treated as a \textit{static shape},
not a sequence to continue: the CNN pools over all 30 cycles simultaneously to
extract curvature, rate of decline, and early-knee features that characterise
where this cell sits in the manufacturing distribution relative to the
pretrained population.

\subsubsection{Phase 3: Fine tuning and MC-Dropout}
The pretrained LSTM is fine tuned
on the first \textit{train\_pct}\% of the target cell's SOH trajectory,
starting from the cell-specific initial state $(h_0, c_0)$ produced by the
fingerprint encoder.
The fingerprint encoder weights, learned in Phase~1, are held \textit{frozen}
throughout fine tuning; only the LSTM and output-layer weights are updated.
Training and inference use fixed-length look-back windows of $W=20$ cycles drawn
from the target cell's SOH. Each window is processed independently: the LSTM is
re-initialised to the fingerprint-derived state $(h_0, c_0)$ at the first step of
every window rather than carrying hidden state across windows, so each prediction
is conditioned on both the cell-level fingerprint and its local $W$-cycle history.
The fine tuning objective is:
\begin{equation}
  \mathcal{L}_{\mathrm{FT}} = \frac{1}{K}\sum_{t=1}^{K}\!\left(\hat{y}_t - y_t\right)^2,
  \quad K = \left\lfloor \frac{p}{100} \cdot T \right\rfloor
  \label{eq:ft_loss}
\end{equation}
where $y_t$ is the observed SOH at cycle $t$, $\hat{y}_t$ is the model
prediction, $T$ is the total cycle count of the target cell, and $p \in
\{30, 50, 70\}$.
Dropout (rate 0.2) applied to the LSTM output during both training and
inference provides uncertainty estimates via 50 stochastic forward passes.
The 95\% prediction interval is $\hat{y} \pm 1.96\,\hat{\sigma}$ (an approximate
MC-Dropout interval, not a calibrated confidence interval), where $\hat{\sigma}$
is the empirical standard deviation across the 50 passes.

\subsection{Pretraining-Run Variance and Checkpoint Selection}\label{sec:selection}

Pretraining a neural network is stochastic: independent runs on the identical
corpus, differing only in random initialisation, converge to representations of
very different transfer quality.
Across five deterministic pretraining runs of Phase~1, downstream SIT
improvement over scratch ranged from $-32\%$ to $+52\%$: three runs learned
genuine degradation dynamics while two collapsed toward a persistence
solution that transfers poorly.
Crucially, \emph{source-domain criteria failed to identify the collapsed
runs}: MIT validation loss was essentially identical for good and collapsed
runs (${\approx}3.4\times10^{-3}$), one-step-ahead RMSE on held-out MIT cells
was \emph{inverted} (collapsed runs win it, because persistence excels at
one-step prediction), and recursive multi-step rollout was uninformative.

We therefore adopt a simple target-domain selection protocol:
\begin{enumerate}
  \item pretrain $k=5$ candidate checkpoints with independent seeds;
  \item evaluate each candidate's fine-tuned RMSE on the three designated
        \emph{calibration cells} (Section~\ref{sec:eval});
  \item deploy the checkpoint with the lowest mean calibration RMSE.
\end{enumerate}
The calibration cells are chosen by an a-priori rule (the numerically first
cell of each condition group) and are excluded from all reported results.
The protocol is robust to the choice of calibration cells: exhaustively
enumerating all $\binom{20}{3}=1140$ possible three-cell calibration sets,
95.5\% select a checkpoint with positive test improvement.
In deployment terms the protocol is realistic: a manufacturer qualifying a
battery product retains laboratory cells of that product, and three such
cells suffice.
All results in Section~\ref{sec:results} use the selected checkpoint; the
full five-run sensitivity study is released with the code.

\subsection{Mechanism Contrast: Fingerprint Encoder vs.\ Fine Tuning}\label{sec:mechanism}

Although Phases 2 and 3 both draw on the same early-cycle SOH data, they
use it in fundamentally different ways, a distinction that directly motivates
the ablation study in Section~\ref{sec:ablation_design}.
Table~\ref{tab:mechanism_contrast} makes this contrast explicit.

\begin{table*}[!t]
\caption{Comparison of the two mechanisms that use early-cycle SOH data.}
\label{tab:mechanism_contrast}
\centering
\begin{tabularx}{\textwidth}{l X X}
\toprule
 & \textbf{Incremental fine tuning} & \textbf{Fingerprint encoder} \\
\midrule
\textit{Input view} & Time series $s_1 \to s_2 \to \cdots \to s_K$ & Static shape $\mathbf{s} \in \mathbb{R}^{30}$ \\
\textit{Operation}  & Adjust LSTM weights $\theta$ to minimise $\mathcal{L}_{\mathrm{FT}}$ & Map shape to $(h_0, c_0)$ via CNN \\
\textit{Output}     & Updated model weights & Cell-specific initial state \\
\textit{Question answered} & What comes next in this sequence? & What type of degrader is this cell? \\
\textit{Limitation} & Ambiguous for extreme cells in the first $p$\% & Requires a diverse pretraining corpus \\
\bottomrule
\end{tabularx}
\end{table*}

For cells at the extremes of the manufacturing distribution
(catastrophically fast: cell 001-8; exceptionally robust: cells 001-5
and 001-7), the first 50\% of the SOH trajectory is \textit{ambiguous}:
it resembles a moderate degrader.
Fine tuning on this ambiguous first half fits a moderate trajectory and
projects it forward, substantially underestimating the second half.
The fingerprint encoder, by contrast, conditions the LSTM on the early-life
trajectory shape (curvature and rate of decline across the first 30 cycles),
placing the prediction in the context of the pretrained population.
Whether this conditioning encodes precise cell \emph{identity} or a more
generic in-distribution prior is an empirical question; a shuffled-fingerprint
control in Section~\ref{sec:results} shows the latter dominates, and we
interpret the mechanism accordingly.

\subsection{Ablation Study Design}\label{sec:ablation_design}

To isolate the contribution of each hybrid component, four configurations
are compared:
\begin{enumerate}
  \item \textbf{Scratch LSTM}: Train the hybrid's LSTM backbone
        ($H{=}64$, $W{=}20$) from scratch on the target cell's training split
        only, with no pretraining, no fingerprint, and no differential-evolution
        hyperparameter search. This is the ablation control for the hybrid and is
        distinct from the DE-optimised DE-LSTM baseline of
        Section~\ref{sec:baseline}; the two therefore differ in RMSE.
  \item \textbf{Zero-shot (pretrain + fingerprint, no fine tuning)}: the MIT
        pretrained model with the fingerprint-derived initial state applied
        directly to the target cell, with \textit{no fine tuning} on target
        data (here \textit{zero-shot} denotes the absence of target-cell fine
        tuning; the first 30 cycles are still used to compute the fingerprint).
        Isolates the value of pretraining and the fingerprint before any
        target-cell adaptation.
  \item \textbf{Pretrain + fine tune (no fingerprint)}: MIT pretrained
        model fine tuned on the target cell's training split, but without
        the fingerprint encoder.
        Quantifies the marginal value of the fingerprint over fine tuning alone.
  \item \textbf{Full hybrid}: All three phases (pretrain + fingerprint + fine-tune).
\end{enumerate}

\subsection{Reproducibility and Determinism}\label{sec:repro}

All experiments run under a fully deterministic configuration: a single
global seed set via \texttt{tf.keras.utils.set\_random\_seed} (which, unlike
\texttt{tf.random.set\_seed} alone, also seeds Keras weight initialisers),
TensorFlow op-determinism enabled, and float32 precision throughout.
Under this configuration every number reported in this paper reproduces
byte-for-byte from a fresh clone of the released code.
The repository includes a single-command verification script
(\texttt{verify\_determinism.sh}) that exercises each experiment family's
actual training and inference code path twice in separate processes and
byte-compares the outputs, so that reviewers and users can confirm
determinism on their own hardware before relying on the pipeline.
We adopted this discipline after observing that library-default (unseeded)
training produced run-to-run variation large enough to alter experimental
conclusions (Sections~\ref{sec:baseline} and~\ref{sec:brittleness}); we
regard it as a necessary condition for benchmark claims in this field.
